\subsection{Nonreduced bases: two opposite examples}

\begin{proposition}[C-RD3: nilpotent invariance when $p$ is locally nilpotent]\label{prop:C-RD3}
\status{Proved.}
Let $S$ be locally noetherian, suppose $p$ is locally nilpotent on $S$, and
let $U\subset S$ be open. Write $S_0=S_{\mathrm{red}}$ and
$U_0=U\times_SS_0$. A triple $(A_U,M,\varphi)$ lies in the essential image
of $F_S$ if and only if its reduction lies in the essential image of
$F_{S_0}$. If $F_{S_0}$ is fully faithful, so is $F_S$. In particular,
if $F_{S_0}$ is an equivalence, then $F_S$ is an equivalence.
\end{proposition}
\begin{proof}
The forward implication is restriction. Conversely let $A_0/S_0$
realize the reduced triple, with the given isomorphism
$A_0[p^\infty]\simeq M|_{S_0}$. Locally on the noetherian scheme $S$,
the nilradical is nilpotent. The Serre--Tate equivalence
\cite[Theorem~1.2.1]{Katz1981SerreTate} constructs an abelian scheme
$A/S$ lifting $A_0$ and realizing $M$. Its restriction to $U$ and $A_U$
have identical reduction and identical $p$-divisible group, with the
specified identification. The full faithfulness in the same theorem
provides their unique compatible isomorphism. These local constructions
glue by that uniqueness. For morphisms, first descend the reduced pair
using full faithfulness of $F_{S_0}$ and then apply full faithfulness of
Serre--Tate. This proves every assertion.
\end{proof}

\begin{corollary}[C-RD4: the proposed reducedness criterion is false]\label{cor:C-RD4}
\status{Refuted, with a positive counterexample.}
For $k=\overline{\mathbf F}_p$, let
\[
 S=\operatorname{Spec}k[[t]][\epsilon]/(\epsilon^2),\qquad U=D(t).
\]
The hypotheses a) and b) of 100:3 hold, the local ring at the point of
$T=S\setminus U$ is nonreduced, and nevertheless $F_S$ is an equivalence
in all relative dimensions.
\end{corollary}
\begin{proof}
The finite surjective map $S_0=\operatorname{Spec}k[[t]]\to S$ supplies
a), since $S_0$ is normal and geometrically unibranch. The etale topoi
are unaffected by the nilpotent thickening. A strictly henselian trait
has connected punctured spectrum, so its local etale cohomology with
locally constant coefficients vanishes in degrees zero and one. This
is $\operatorname{prof}^{\mathrm{et}}_T(S)\geq2$; it verifies b).

Over the trait $S_0$, the existence of an extension of $A_{U_0}[p^\infty]$
implies good reduction of $A_{U_0}$ by de Jong's good-reduction criterion
\cite[Theorem~3, citing Theorem~2.5 of the original paper]{deJong1998ICM}.
The given identification with $M_0$ extends uniquely, by the extension
theorem for homomorphisms of Barsotti--Tate groups
\cite[Theorem~2]{deJong1998ICM}. Morphisms of the abelian schemes over a
trait extend by the Neron mapping property. Thus $F_{S_0}$ is an
equivalence. Proposition~\ref{prop:C-RD3} gives the assertion for $S$.
The nonzero element $\epsilon$ proves nonreducedness at $T$.
\end{proof}

\begin{proposition}[C-RD5: an explicit failure of 100:3(2)]\label{prop:C-RD5}
\status{Proved counterexample; R4 rechecked.}
Let
\[
 S=\operatorname{Spec}\mathbf Z_5[\epsilon]/(\epsilon^2),\quad U=D(5),
 \quad E/\mathbf Z_5:\ y^2=x^3+x+1.
\]
Put $M=E[5^\infty]\times_{\mathbf Z_5}S$, and over $U$ take
\[
 A_U:\quad y^2=x^3+x+(1+\epsilon).
\]
There is a unique isomorphism $M|_U\simeq A_U[5^\infty]$ reducing to the
identity modulo $\epsilon$. This triple does not extend to an abelian
scheme over $S$, although a) and b) of 100:3 hold.
\end{proposition}
\begin{proof}
Here $\Delta(E)=-496$ is a $5$-adic unit, so $E$ is an elliptic scheme.
The two generic $5$-divisible groups are etale; invariance of finite
etale schemes under nilpotent thickenings gives the stated compatible
isomorphisms at all levels. Hypothesis a) holds via
$\operatorname{Spec}\mathbf Z_5\to S$, and b) follows from the connected
punctured strict localization of this trait, as above.

If $A/S$ realized the triple, its reduction modulo $\epsilon$ would be
$E$: both are good models of the same elliptic curve over $\mathbf Q_5$.
For every $n\geq1$, apply Serre--Tate over
$(\mathbf Z_5/5^n)[\epsilon]/(\epsilon^2)$. The specified constant
$5$-divisible-group deformation forces $A$ modulo $5^n$ to be the
constant deformation of $E$ modulo $5^n$. Hence
$j(A)-j(E)$ lies in $\bigcap_n5^n\mathbf Z_5[\epsilon]=0$.
On the other hand, the generic fibre prescribed above has
\[
 j(A_U)=\frac{6912}{31+54\epsilon}
 =\frac{6912}{31}-\frac{373248}{961}\epsilon,
\]
whose $\epsilon$ coefficient is nonzero. This contradiction proves the
claim. It also identifies exactly the failure of the proposed gluing
at 100:7: equality of the generic etale $p$-divisible groups supplies no
isomorphism between these two characteristic-zero infinitesimal
abelian-scheme deformations.
\end{proof}
